\documentclass[12pt]{article}
\usepackage[spanish]{babel}
\usepackage[a4paper, margin=2.5cm]{geometry}
\usepackage{graphicx}
\usepackage{float}
\usepackage{textpos}
\usepackage{tikz}
\usetikzlibrary{arrows.meta}
\usepackage{listings}
\usepackage{xcolor}
\usepackage{amsmath}
\usepackage{fancyhdr}
\setlength{\headheight}{2.4cm}
\pagestyle{fancy}
\fancyhf{}
\lhead{\includegraphics[height=1.8cm]{head.png}}
\cfoot{\thepage}
\renewcommand{\headrulewidth}{0pt}

\lstset{
    basicstyle=\ttfamily\small,
    keywordstyle=\color{blue}\bfseries,
    commentstyle=\color{gray},
    stringstyle=\color{orange},
    numbers=left,
    numberstyle=\tiny\color{gray},
    frame=single,
    breaklines=true,
    showstringspaces=false
}
\newcommand{\cmd}[1]{\colorbox{gray!20}{\texttt{#1}}}
\newcommand{\nota}[1]{\textit{\textcolor{gray}{#1}}}
\newcommand{\imagen}[3]{%
  \begin{figure}[H]
    \centering
    {\setlength{\fboxrule}{1pt}\setlength{\fboxsep}{4pt}%
     \fcolorbox{black!60}{white}{%
       \setlength{\fboxrule}{0.5pt}\setlength{\fboxsep}{8pt}%
       \fcolorbox{black!35}{white}{\includegraphics[width=#2\textwidth]{#1}}}}
    \caption{#3}
    \label{fig:#1}
  \end{figure}
}

\begin{document}

% Portada
\thispagestyle{empty}

\begin{tikzpicture}[remember picture, overlay]
    % Logo y información de la universidad en la parte superior
    \node[anchor=north west] at (current page.north west) {%
        \begin{minipage}{\textwidth}
            \vspace{1cm}
            \hspace{0.5cm}
            \begin{minipage}{2.5cm}
                \includegraphics[width=2.5cm]{Utalca.png}
            \end{minipage}
            \hspace{0.5cm}
            \begin{minipage}{10cm}
                \raggedright
                {\small\textbf{UNIVERSIDAD DE TALCA}} \\
                {\small\textbf{FACULTAD DE INGENIERÍA}} \\
                {\small\textbf{DEPARTAMENTO DE CIENCIAS DE LA COMPUTACIÓN}}
            \end{minipage}
        \end{minipage}
    };
    
    % Información en la parte inferior derecha
    \node[anchor=south west] at ([xshift=2cm, yshift=1cm]current page.south west) {%
        \begin{tabular}{@{}ll@{}}
            \textbf{Fecha:}  & 11/11/2026 \\[0.3cm]
            \textbf{Autores:}   & Gaston Villanueva \\
                                & Nicolas Salas \\
                                & Krisstal  Hernandez \\[0.3cm]    
            \textbf{e-mail:} & \texttt{gvillanueva19@alumnos.utalca.cl} \\ 
                            &  \texttt{nsalas21@alumnos.utalca.cl} \\
                            & \texttt{khernandez22@alumnos.utalca.cl} \\ [0.3cm]
            \textbf{Profesor:} & Renzo Angles \\[0.3cm]
        \end{tabular}
    };
\end{tikzpicture}

% Espacio en blanco
\vspace*{4cm}

% Título y subtítulo centrados
\begin{center}
    \vspace*{2cm}
    
    {\huge\textbf{Informe Proyecto Grupal de Unidad }}

    {\vspace{1cm}}

    {\Large\textbf{Grupo 5}}
    
    {\vspace{1cm}}
    
    {\Large Módulo de EXTRACCION Y GESTION DE DATOS MASIVOS}

    {\vspace{1cm}}

    {\Large Etapa 1}
    
\end{center}

\newpage
%\input{indice}
%\input{indice.tex}
%\input{introduccion}
%\input{validacionRut}
%\input{triggers}
%\input{procedimientosAlmacenados}
%\input{vistas}
%\input{reglasAvanzadas}
%\input{resultadosPruebas.tex}
%\input{conclusion}

\newpage
% ============================================================
\section{Propuesta}
\subsection{Herramienta elegida.}
\subsubsection*{Apache Cassandra} La herramienta corresponde a una base de datos NoSQL distribuida de tipo
\textit{wide-column}, sin nodo maestro (arquitectura \textit{peer-to-peer}), con replicación configurable
y consistencia ajustable por consulta. 


\subsection{¿Por qué se eligió esta base de datos?}
\begin{itemize}\itemsep0pt
  \item \textbf{Encaja con el tipo de dato:} las lecturas de sensores son series de tiempo de alta tasa de
        escritura, el caso de uso para el que Cassandra está diseñada (escritura mediante \textit{commit log},
        \textit{memtable} y \textit{SSTable}).
  \item \textbf{Complementa el ecosistema visto:} a diferencia de HDFS/Hive, que están orientados a lotes y
        análisis, Cassandra atiende lecturas y escrituras de baja latencia por clave y no tiene punto único de falla.
  \item \textbf{Operable localmente:} se instala con Docker Compose; el clúster de 3 nodos levantó y operó
        sin problemas en el equipo del grupo.
  \item \textbf{Permite experimentos claros para la Etapa 3:} consistencia (\texttt{ONE} vs.\ \texttt{QUORUM}),
        tolerancia a fallos al detener nodos y escalabilidad al variar el número de nodos.
\end{itemize}

\subsection{Dataset tentativo: Smart Manufacturing IoT-Cloud Monitoring Dataset}
\noindent
\textbf{Fuente:} Kaggle \\
\textbf{Autor:} \texttt{ziya07} \\
\textbf{Tipo de dominio:} Público \\
\textbf{Volumen:} 100\,000 lecturas (1 por minuto) \\
\textbf{Máquinas:} 50 máquinas industriales \\
\textbf{Periodo:} 1 de enero -- 11 de marzo de 2025 \\
\textbf{Variables:} 13 columnas (temperatura, vibración, humedad, presión, consumo de energía, estado de la máquina, indicador de anomalía, vida remanente estimada, tipo de falla, riesgo de caída y mantención requerida).

\vspace{0.6em}
Este dataset está orientado al \textbf{mantenimiento predictivo} y a la \textbf{detección de anomalías}. 
Dado que el autor no precisa su procedencia, el grupo lo trata como un conjunto de datos de sensores 
industriales, sin afirmar que provenga de una planta real. Para la Etapa 3, el volumen puede ampliarse 
replicando el conjunto sobre un mayor número de máquinas y un horizonte temporal más extenso.
\subsection{Herramientas probadas y descartadas.}
\textbf{Ninguna}. Se eligió como primera opción Cassandra y las
pruebas básicas resultaron sin necesidad de evaluar otra herramienta; por eso no se declaran alternativas
descartadas.
\newpage
\section{Pruebas básicas}
\subsection{Modelo de datos (\texttt{schema.cql})}
A continuación se mostrará de manera visual el esquema del Keyspace \cmd{manufactura\_iot} con \cmd{NetworkTopologyStrategy} y factor de replicación 3.
Se usa una tabla por consulta; el campo \texttt{mes} acota el tamaño de cada partición.

\begin{figure}[h!]
\centering
\newcommand{\cqltable}[4]{% x, y, titulo, filas
  \node[anchor=north west, fill=black!75, text=white, minimum width=10.5cm, minimum height=0.7cm,
        font=\bfseries\small\ttfamily, draw=black] at (#1,#2) {#3};
  \foreach \n/\t/\k [count=\i] in {#4} {
    \pgfmathsetmacro{\yy}{#2-0.7-(\i-1)*0.5}
    \if\k P\def\fillc{blue!20}\def\kt{PK}\else
      \if\k C\def\fillc{orange!30}\def\kt{CK}\else\def\fillc{white}\def\kt{}\fi\fi
    \node[anchor=north west, draw, fill=\fillc, minimum width=5.5cm, minimum height=0.5cm,
          font=\footnotesize\ttfamily, align=left, inner sep=4pt] at (#1,\yy) {\n};
    \node[anchor=north west, draw, fill=\fillc, minimum width=3cm, minimum height=0.5cm,
          font=\footnotesize\ttfamily] at (#1+5.5,\yy) {\t};
    \node[anchor=north west, draw, fill=\fillc, minimum width=2cm, minimum height=0.5cm,
          font=\footnotesize\bfseries] at (#1+8.5,\yy) {\kt};
  }
}
\begin{tikzpicture}[x=1cm,y=1cm]
% Tabla 1
\cqltable{0}{0}{lecturas\_por\_maquina}{%
  machine\_id/int/P, mes/text/P, ts/timestamp/C, temperature/float/N, vibration/float/N,
  humidity/float/N, pressure/float/N, energy\_consumption/float/N, machine\_status/int/N,
  anomaly\_flag/int/N, predicted\_remaining\_life/int/N, failure\_type/text/N,
  downtime\_risk/float/N, maintenance\_required/int/N}
\node[anchor=north west, font=\scriptsize] at (0,-7.9) {\texttt{PRIMARY KEY ((machine\_id, mes), ts)} \quad \texttt{CLUSTERING ORDER BY (ts DESC)}};
% Flecha
\draw[-{Stealth[length=3mm]}, thick] (5.25,-8.4) -- node[right, font=\scriptsize, align=left]{\ desnormalización: solo lecturas con falla} (5.25,-9.4);
% Tabla 2
\cqltable{0}{-9.5}{anomalias\_por\_tipo}{%
  failure\_type/text/P, mes/text/P, ts/timestamp/C, machine\_id/int/C, temperature/float/N,
  vibration/float/N, pressure/float/N, anomaly\_flag/int/N, predicted\_remaining\_life/int/N,
  downtime\_risk/float/N}
\node[anchor=north west, font=\scriptsize] at (0,-15.7) {\texttt{PRIMARY KEY ((failure\_type, mes), ts, machine\_id)} \quad \texttt{ORDER BY (ts DESC, machine\_id ASC)}};
% Leyenda
\node[draw, fill=blue!20, minimum width=0.5cm, minimum height=0.35cm] at (0.4,-16.7) {};
\node[anchor=west, font=\scriptsize] at (0.8,-16.7) {Clave de partición};
\node[draw, fill=orange!30, minimum width=0.5cm, minimum height=0.35cm] at (4.2,-16.7) {};
\node[anchor=west, font=\scriptsize] at (4.6,-16.7) {Clave de clustering};
\node[draw, fill=white, minimum width=0.5cm, minimum height=0.35cm] at (8,-16.7) {};
\node[anchor=west, font=\scriptsize] at (8.4,-16.7) {Columna regular};
\end{tikzpicture}
\caption{Esquema de las tablas de Cassandra en el keyspace \texttt{manufactura\_iot}.}
\end{figure}

\newpage
\subsection{Instalación reproducible}
El clúster se define en un único archivo \cmd{docker-compose.yml}: tres nodos \texttt{cassandra:5.0}
(\texttt{cass1}, \texttt{cass2}, \texttt{cass3}) en un mismo datacenter y tres \textit{racks}, con
\cmd{GossipingPropertyFileSnitch} y \cmd{heap} de \cmd{512\,MB} por nodo. Los nodos arrancan en cadena
(\cmd{depends\_on} con \cmd{healthcheck}), de modo que cada uno se une al anillo cuando el anterior
está listo. Todo se levanta con un comando:

\begin{lstlisting}[language=bash]
docker compose up -d #referirse al archivo docker-compose.yml para ver su contenido
\end{lstlisting}
\subsubsection{Verificación del anillo} Asi con \cmd{nodetool status}: los tres nodos aparecen en estado \cmd{UN}
(\cmd{Up/Normal}), cada uno en su \textit{rack} y con 100\,\% de propiedad efectiva, coherente con un
factor de replicación 3 sobre 3 nodos.

\begin{lstlisting}
$ docker exec cass1 nodetool status
--  Address     Load        Tokens  Owns (effective)  Rack
UN  172.20.0.3  85.1 KiB    16      100.0%            rack2
UN  172.20.0.2  3.55 MiB    16      100.0%            rack1
UN  172.20.0.4  119.69 KiB  16      100.0%            rack3
\end{lstlisting}

\subsection{Carga del dataset}


\subsubsection{Preparación.} El CSV original (\cmd{smart\_manufacturing\_data.csv}, 100\,000 filas) se
transforma con \cmd{preparar\_datos.py}: se renombra \texttt{timestamp} a \texttt{ts} (palabra clave de
CQL) y se agrega la columna \texttt{mes}, que se usa como parte de la clave de partición. El script genera
\cmd{lecturas.csv} (100\,000 filas) y \cmd{anomalias.csv} (8\,101 filas con falla).

\subsection{Modelo de datos}
\subsubsection{Se diseñó una tabla por consulta:}
Se diseñó una tabla por cada tipo de consulta, optimizando el acceso según el patrón de uso:

\begin{itemize}\itemsep0.4em
  \item \textbf{\texttt{lecturas\_por\_maquina}}  
  Clave primaria: \texttt{((machine\_id, mes), ts)}  
  Consulta principal: obtener las lecturas de una máquina en un mes determinado, ordenadas de la más reciente a la más antigua.  
  Se particiona por mes para mantener un tamaño de partición acotado (máximo 981 filas en este dataset).

  \item \textbf{\texttt{anomalias\_por\_tipo}}  
  Clave primaria: \texttt{((failure\_type, mes), ts, machine\_id)}  
  Consulta principal: recuperar todas las fallas de un tipo específico ocurridas en un mes, provenientes de cualquier máquina.
\end{itemize}

\nota{El esquema completo se encuentra definido en el archivo \texttt{schema.cql}.}

\newpage
\subsection{Carga}
Los datos se cargaron mediante el comando \texttt{COPY} de \texttt{cqlsh}, utilizando nivel de consistencia \texttt{QUORUM}:

\begin{lstlisting}[language=bash,caption={Comandos de carga}]
CONSISTENCY QUORUM;

COPY manufactura_iot.lecturas_por_maquina 
  (machine_id, mes, ts, temperature, vibration, humidity, pressure, 
   energy_consumption, machine_status, anomaly_flag, 
   predicted_remaining_life, failure_type, downtime_risk, 
   maintenance_required) 
FROM '/datos/lecturas.csv' WITH HEADER = true;

COPY manufactura_iot.anomalias_por_tipo 
  (failure_type, mes, ts, machine_id, temperature, vibration, 
   pressure, anomaly_flag, predicted_remaining_life, downtime_risk) 
FROM '/datos/anomalias.csv' WITH HEADER = true;
\end{lstlisting}

Posteriormente, se verificó la correcta inserción de los datos:

\begin{lstlisting}
cqlsh> SELECT count(*) FROM manufactura_iot.lecturas_por_maquina;
 count
--------
 100000

cqlsh> SELECT count(*) FROM manufactura_iot.anomalias_por_tipo;
 count
-------
   8101
\end{lstlisting}
{\vspace{1cm}}
Ambos conteos coinciden con el número de filas de los archivos de origen. \par
{\vspace{1cm}}
\nota{\texttt{cqlsh} advierte que la agregación \texttt{count(*)} se realizó sin \textit{partition key}. 
Esta advertencia es aceptable en este contexto, ya que se trata únicamente de una verificación puntual.}
\newpage

\subsection{Operaciones representativas}
Las tres operaciones se ejecutan con \cmd{operaciones.cql}, con consistencia \texttt{QUORUM}:
\begin{lstlisting}[language=bash]
docker exec -it cass1 cqlsh -f /etc/cassandra-lab/operaciones.cql
\end{lstlisting}
\subsubsection{Escribir}
Inserción de una nueva lectura de la máquina 15 en ambas tablas (una por consulta):
\begin{lstlisting}[language=SQL]
INSERT INTO lecturas_por_maquina (machine_id, mes, ts, temperature, vibration, humidity, pressure,
  energy_consumption, machine_status, anomaly_flag, predicted_remaining_life, failure_type,
  downtime_risk, maintenance_required)
VALUES (15, '2025-03', '2025-03-11 10:40:00+0000', 99.5, 45.0, 50.0, 3.0, 4.2, 2, 1, 12,
  'Overheating', 1.0, 1);
INSERT INTO anomalias_por_tipo (failure_type, mes, ts, machine_id, temperature, vibration, pressure,
  anomaly_flag, predicted_remaining_life, downtime_risk)
VALUES ('Overheating', '2025-03', '2025-03-11 10:40:00+0000', 15, 99.5, 45.0, 3.0, 1, 12, 1.0);
\end{lstlisting}

\noindent \texttt{cqlsh} no imprime salida cuando una escritura se completa; la verificación está en la consulta siguiente, donde la fila
insertada (10:40) aparece en ambas tablas.

\subsubsection{Consultar}
Últimas cinco lecturas de la máquina 15 en marzo, y últimas cinco fallas por sobrecalentamiento del mes.
Ambas consultas filtran por la clave de partición completa, por lo que Cassandra lee una sola partición y
los datos ya están ordenados por \texttt{ts} descendente:

\begin{lstlisting}[language=SQL]
SELECT ts, temperature, vibration, pressure, failure_type FROM lecturas_por_maquina
WHERE machine_id = 15 AND mes = '2025-03' LIMIT 5;

SELECT ts, machine_id, temperature, predicted_remaining_life FROM anomalias_por_tipo
WHERE failure_type = 'Overheating' AND mes = '2025-03' LIMIT 5;
\end{lstlisting}
\newpage
\subsection{Salidas de la terminal}
\imagen{imagenesCapturasEvidencias/captura_terminal.png}{0.85}{Ejecución del archivo \texttt{operaciones.cql} mostrando las últimas lecturas de la máquina 15.}
\imagen{imagenesCapturasEvidencias/captura_terminal2.png}{0.85}{Ejecución del archivo \texttt{operaciones.cql}}\par
La fila insertada en la operación anterior aparece primera en ambas consultas, lo que confirma que la
escritura quedó replicada con \texttt{QUORUM} y que el orden descendente por \texttt{ts} se resuelve
directamente desde la partición.

\newpage
\subsection{Agregar}
Resumen de la máquina 15 en febrero con funciones de agregación nativas, y conteo de fallas por vibración:
\begin{lstlisting}[language=SQL]
SELECT count(*) AS lecturas, avg(temperature) AS temp_prom, max(vibration) AS vib_max,
       sum(anomaly_flag) AS anomalias, min(predicted_remaining_life) AS vida_min
FROM lecturas_por_maquina WHERE machine_id = 15 AND mes = '2025-02';

SELECT count(*) AS fallas_vibracion FROM anomalias_por_tipo
WHERE failure_type = 'Vibration Issue' AND mes = '2025-02';
\end{lstlisting}
\imagen{imagenesCapturasEvidencias/conteoVib.png}{0.85}{Salida de terminal con las vibraciones \texttt{operaciones.cql}}\par
\imagen{imagenesCapturasEvidencias/fallasVib.png}{0.2}{Fallas de vibraciones \texttt{operaciones.cql}}\par

La máquina 15 registró 775 lecturas en febrero, con 68 marcadas como anomalía, y en el mismo mes hubo 1\,239
fallas por vibración en el conjunto de las 50 máquinas.



\end{document}